\begin{derivation}[从\eqnrefs{eq:qed-reduced-exchange-r13}到\eqnrefs{eq:coulomb-Vq-r13}]
正文中的非相对论电流展开不能只作为一个“已知低速结论”代入；它直接来自前文已经固定归一化的正能 Dirac 自旋子。对质量为 $m$ 的外线，记
\begin{equation*}
 D_{\mathbf p}\equiv \frac{E(\mathbf p)}{c}+mc,
 \qquad
 N_{\mathbf p}\equiv\sqrt{D_{\mathbf p}},
\end{equation*}
则正文\eqnrefs{eq:explicit-u} 可写成
\begin{equation*}
 u_s(p)=N_{\mathbf p}
 \begin{pmatrix}
 \xi_s\\[1mm]
 \dfrac{\boldsymbol\sigma_{\rm P}\cdot\mathbf p}{D_{\mathbf p}}\xi_s
 \end{pmatrix},
 \qquad
 \xi_{s'}^\dagger\xi_s=\delta_{s's}.
\end{equation*}
由精确色散
$E(\mathbf p)=\sqrt{m^2c^4+c^2\mathbf p^2}$，在
$|\mathbf p|/(mc)\ll1$ 时
\begin{align*}
 \frac{E(\mathbf p)}{c}
 &=mc+\frac{\mathbf p^2}{2mc}
 +O\!\left(\frac{|\mathbf p|^4}{m^3c^3}\right),\\
 D_{\mathbf p}
 &=2mc+\frac{\mathbf p^2}{2mc}
 +O\!\left(\frac{|\mathbf p|^4}{m^3c^3}\right),\\
 N_{\mathbf p}
 &=\sqrt{2mc}\left[
 1+\frac{\mathbf p^2}{8m^2c^2}
 +O\!\left(\frac{|\mathbf p|^4}{m^4c^4}\right)
 \right].
\end{align*}

时间分量的 Dirac 电流为
\begin{align*}
 \bar u_{s'}(p')\gamma^0u_s(p)
 &=u_{s'}^\dagger(p')u_s(p)\\
 &=N_{\mathbf p'}N_{\mathbf p}\,
 \xi_{s'}^\dagger\left[
 I_2+
 \frac{(\boldsymbol\sigma_{\rm P}\cdot\mathbf p')
       (\boldsymbol\sigma_{\rm P}\cdot\mathbf p)}
 {D_{\mathbf p'}D_{\mathbf p}}
 \right]\xi_s.
\end{align*}
第二项已经含两个小动量；同时
$N_{\mathbf p'}N_{\mathbf p}=2mc+O(|\mathbf p|^2/(mc),|\mathbf p'|^2/(mc))$，故
\begin{equation*}
 \boxed{
 \bar u_{s'}(p')\gamma^0u_s(p)
 =2mc\,\delta_{s's}
 +O\!\left(\frac{|\mathbf p|^2+|\mathbf p'|^2}{mc}\right).}
\end{equation*}

空间分量不能仅靠量纲判断。由
$\bar u\gamma^iu=u^\dagger\alpha^iu$ 与
$\alpha^i=\begin{psmallmatrix}0&\sigma_{\rm P}^i\\
\sigma_{\rm P}^i&0\end{psmallmatrix}$，有
\begin{align*}
 \bar u_{s'}(p')\gamma^iu_s(p)
 ={}&N_{\mathbf p'}N_{\mathbf p}\,
 \xi_{s'}^\dagger\left[
 \frac{\sigma_{\rm P}^i(\boldsymbol\sigma_{\rm P}\cdot\mathbf p)}{D_{\mathbf p}}
 +\frac{(\boldsymbol\sigma_{\rm P}\cdot\mathbf p')\sigma_{\rm P}^i}{D_{\mathbf p'}}
 \right]\xi_s.
\end{align*}
利用
$\sigma_{\rm P}^i\sigma_{\rm P}^j
=\delta^{ij}I_2+\ii\epsilon^{ijk}\sigma_{\rm P}^k$，领先阶成为
\begin{equation*}
 \boxed{
 \bar u_{s'}(p')\gamma^iu_s(p)
 =\xi_{s'}^\dagger\left[
 (p+p')^i I_2
 +\ii\epsilon^{ijk}(p-p')^j\sigma_{\rm P}^k
 \right]\xi_s
 +O\!\left(\frac{|\mathbf p|^3+|\mathbf p'|^3}{m^2c^2}\right).}
\end{equation*}
所以时间分量为 $O(mc)$，空间分量为 $O(|\mathbf p|)$；这才是静态极限中电荷密度压倒三维电流的数学原因。

对两条外费米子线分别应用这一结果，并取不翻转自旋通道，正文\eqnrefs{eq:qed-reduced-exchange-r13} 给出
\begin{align*}
 \mathcal M_\gamma
 &=\frac{g_1g_2}{q^2}
 J_1^\mu J_{2\mu}\\
 &=\frac{g_1g_2}{q^2}
 \left[J_1^0J_2^0-\mathbf J_1\cdot\mathbf J_2\right]\\
 &=\frac{4m_1m_2c^2g_1g_2}{q^2}
 +O\!\left(\frac{|\mathbf p|^2}{m_1m_2c^2}\right).
\end{align*}
近静态交换再给
\begin{equation*}
 q^2=(q^0)^2-\mathbf q^2
 =-\mathbf q^2\left[
 1-\frac{(q^0)^2}{\mathbf q^2}
 \right],
\end{equation*}
故
\begin{equation*}
 \mathcal M_\gamma^{\rm NR}(\mathbf q)
 =-\frac{4m_1m_2c^2g_1g_2}{\mathbf q^2}
 +O\!\left(
 \frac{|\mathbf p|^2}{\mathbf q^2},
 \frac{m_1m_2c^2g_1g_2(q^0)^2}{|\mathbf q|^4}
 \right).
\end{equation*}
把领先项代入前文已经从相对论与 Schrödinger 截面归一化推得的匹配式
\eqnrefs{eq:relativistic-to-nr-potential-matching-dp75}：
\begin{align*}
 \widetilde V_C(\mathbf q)
 &=-\frac{\hbar^3}{4m_1m_2c}
 \mathcal M_\gamma^{\rm NR}(\mathbf q)\\
 &=\frac{\hbar^3c}{\mathbf q^2}
 \frac{q_1q_2}{\varepsilon_0\hbar c}\\
 &=\boxed{\frac{q_1q_2\hbar^2}
 {\varepsilon_0\mathbf q^2}},
\end{align*}
即正文\eqnrefs{eq:coulomb-Vq-r13}。随后只剩三维 Fourier 变换
\begin{equation*}
 \int\frac{\dd^3q}{(2\pi\hbar)^3}
 \frac{\ee^{\ii\mathbf q\cdot\mathbf r/\hbar}}{\mathbf q^2}
 =\frac{1}{4\pi\hbar^2r},
\end{equation*}
因此
$V_C(r)=q_1q_2/(4\pi\varepsilon_0r)$。这最后一步是标准 Fourier 积分；真正的物理内容在前面的三步：Dirac 电流低速展开、近静态传播子极限以及相对论振幅到势核的归一化匹配。
\end{derivation}



\begin{derivation}[从\eqnrefs{eq:emu-M2-trace}到\eqnrefs{eq:emu-M2-massive-invariants-dp49}]
取外态

\begin{equation*}
 e^-(p_1,s_1)+\mu^-(p_2,s_2)
 \to e^-(p_3,s_3)+\mu^-(p_4,s_4),
\end{equation*}
则 $q=p_1-p_3$、$q^2=t$，振幅为
\begin{align*}
 \ii\mathcal M_t
 &=(-\ii g_{\rm em})\bar u_e(p_3)\gamma^\mu u_e(p_1)
 \frac{-\ii\eta_{\mu\nu}}{t}
 (-\ii g_{\rm em})\bar u_\mu(p_4)\gamma^\nu u_\mu(p_2),\\
 \mathcal M_t
 &=\frac{g_{\rm em}^2}{t}
 [\bar u_3\gamma^\mu u_1]
 [\bar u_4\gamma_\mu u_2].
 \end{align*}
未偏振实验对电子和 $\mu$ 子初态各平均两种自旋，所以
\begin{align*}
 \overline{|\mathcal M_t|^2}
 =&\frac{g_{\rm em}^4}{t^2}\frac14
 \Tr[(\slashed p_3+\me c)\gamma^\mu(\slashed p_1+\me c)\gamma^\nu]\nonumber\\
 &\hspace{2.4cm}\times
 \Tr[(\slashed p_4+m_\mu c)\gamma_\mu(\slashed p_2+m_\mu c)\gamma_\nu].
 \end{align*}
使用 \eqnrefs{eq:qed-basic-tensor-si-r9}，先把两条费米子线的迹分别写成秩二张量：
\begin{align*}
 L_\ee^{\mu\nu}
 &=4\left[
 p_3^\mu p_1^\nu+p_3^\nu p_1^\mu
 -(p_3\!\cdot p_1-\me^2c^2)\eta^{\mu\nu}
 \right],\\
 L_\mu^{\mu\nu}
 &=4\left[
 p_4^\mu p_2^\nu+p_4^\nu p_2^\mu
 -(p_4\!\cdot p_2-m_\mu^2c^2)\eta^{\mu\nu}
 \right].
 \end{align*}
于是 \eqnrefs{eq:emu-M2-trace} 变成
\begin{align*}
 \overline{|\mathcal M_t|^2}
 &=\frac{g_{\rm em}^4}{4t^2}L_\ee^{\mu\nu}L_{\mu,\mu\nu},\\
 &=\frac{g_{\rm em}^4}{t^2}\left\{
 8\left[(p_1\!\cdot p_2)(p_3\!\cdot p_4)
 +(p_1\!\cdot p_4)(p_2\!\cdot p_3)\right]
 +4t(\me^2+m_\mu^2)c^2
 \right\}.
 \end{align*}
这里第二行只做了 Lorentz 指标收缩。再由
\begin{align*}
 2p_1\!\cdot p_2=2p_3\!\cdot p_4
 &=s-(\me^2+m_\mu^2)c^2\equiv S_0,\\
 -2p_1\!\cdot p_4=-2p_2\!\cdot p_3
 &=u-(\me^2+m_\mu^2)c^2\equiv U_0,
\end{align*}
得到
\begin{align*}
 \overline{|\mathcal M_t|^2}
 =\frac{2g_{\rm em}^4}{t^2}
 \left[
 S_0^2+U_0^2+2t(\me^2+m_\mu^2)c^2
 \right],
\end{align*}
即正文\eqnrefs{eq:emu-M2-massive-invariants-dp49}。


\end{derivation}

\begin{derivation}[从\eqnrefs{eq:identical-fermion-exchange-sign-dp68}到\eqnrefs{eq:moller-angular}]
取入射动量为 $p_1,p_2$，出射动量为 $p_3,p_4$。由于两条末态电子不可区分，有两种 Wick 收缩：

\begin{align*}
 \mathcal M_t
 &=\frac{g_{\rm em}^2}{t}
 [\bar u_3\gamma^\mu u_1]
 [\bar u_4\gamma_\mu u_2],\\
 \mathcal M_u
 &=\frac{g_{\rm em}^2}{u}
 [\bar u_4\gamma^\mu u_1]
 [\bar u_3\gamma_\mu u_2].
\end{align*}
交换两个外部费米子算符需要一次奇置换，因此
\begin{equation*}
 {
 \mathcal M_{\mathrm{Moller}}=\mathcal M_t-\mathcal M_u.}
 \end{equation*}
平方以后
\begin{equation*}
 |\mathcal M_t-\mathcal M_u|^2
 =|\mathcal M_t|^2+|\mathcal M_u|^2
 -2\Re(\mathcal M_t\mathcal M_u^*).
\end{equation*}
在 $s\gg \me^2c^2$ 极限，两个平方项就是前面单一交换迹的两种外部标签分配：
\begin{align*}
 \overline{|\mathcal M_t|^2}
 &=2g_{\rm em}^4\frac{s^2+u^2}{t^2},\\
 \overline{|\mathcal M_u|^2}
 &=2g_{\rm em}^4\frac{s^2+t^2}{u^2}.
\end{align*}
干涉项需要单独计算。把四条自旋求和首尾闭合成一个迹，得到
\begin{align*}
 \overline{-2\Re(\mathcal M_t\mathcal M_u^*)}
 &=-\frac{g_{\rm em}^4}{2tu}
 \Tr\!\left[
 \slashed p_3\gamma^\mu\slashed p_1\gamma^\nu
 \slashed p_4\gamma_\mu\slashed p_2\gamma_\nu
 \right].
 \end{align*}
使用前面由 Clifford 代数建立的
\eqnrefs{eq:gamma-two-sandwich-r10,eq:gamma-three-sandwich-r10}，迹依次化为
\begin{align*}
 \mathcal I_{tu}
 &\equiv
 \Tr[\slashed p_3\gamma^\mu\slashed p_1\gamma^\nu
 \slashed p_4\gamma_\mu\slashed p_2\gamma_\nu],\\
 &=-2\Tr[\slashed p_3\gamma^\mu\slashed p_1\slashed p_2
 \gamma_\mu\slashed p_4],\\
 &=-32(p_1\!\cdot p_2)(p_3\!\cdot p_4)
 =-8s^2.
\end{align*}
因此
\begin{equation*}
 \overline{-2\Re(\mathcal M_t\mathcal M_u^*)}
 =4g_{\rm em}^4\frac{s^2}{tu}.
 \end{equation*}
因此
\begin{equation*}
 {
 \overline{|\mathcal M_{\mathrm{Moller}}|^2}
 =2g_{\rm em}^4\left[
 \frac{s^2+u^2}{t^2}
 +\frac{s^2+t^2}{u^2}
 +\frac{2s^2}{tu}
 \right].}
 \end{equation*}
若 $\theta$ 在全 $0\le\theta\le\pi$ 上标记“第 3 条电子”的方向，则 $(p_3,p_4)$ 与 $(p_4,p_3)$ 被重复计数，需要 $1/2!$。使用
\begin{equation*}
 t=-\frac s2(1-\cos\theta),
 \qquad
 u=-\frac s2(1+\cos\theta)
\end{equation*}
可化为
\begin{equation*}
 {
 \frac{\dd\sigma_{\mathrm{Moller}}}{\dd\Omega}
 =\frac{\hbar^2\alpha^2}{2s}
 \frac{(3+\cos^2\theta)^2}{\sin^4\theta}.}
 \end{equation*}
若只在一个不重复的半球区域定义末态，则不再额外除 $2!$，两种约定必须二选一。
\end{derivation}

\begin{derivation}[从\eqnrefs{eq:bhabha-tree-amplitude-r68}到\eqnrefs{eq:bhabha-angular-r68}]
Bhabha 散射同时有湮灭的 $s$ 通道与散射的 $t$ 通道：

\begin{align*}
 \mathcal M_s
 &=\frac{g_{\rm em}^2}{s}
 [\bar v_2\gamma^\mu u_1]
 [\bar u_3\gamma_\mu v_4],\\
 \mathcal M_t
 &=-\frac{g_{\rm em}^2}{t}
 [\bar u_3\gamma^\mu u_1]
 [\bar v_2\gamma_\mu v_4].
 \end{align*}
相对号由把同一个 Wick 收缩次序写成标准外部旋量次序时确定。对无偏振、$s\gg \me^2c^2$，两个平方项分别复用湮灭与交换的迹：
\begin{align*}
 \overline{|\mathcal M_s|^2}
 &=2g_{\rm em}^4\frac{u^2+t^2}{s^2},\\
 \overline{|\mathcal M_t|^2}
 &=2g_{\rm em}^4\frac{u^2+s^2}{t^2}.
\end{align*}
交叉项把四条外部自旋求和s 连成
\begin{align*}
 2\Re\,\overline{\mathcal M_s\mathcal M_t^*}
 &=-\frac{g_{\rm em}^4}{2st}
 \Tr\!\left[
 \slashed p_2\gamma^\mu\slashed p_1\gamma^\nu
 \slashed p_3\gamma_\mu\slashed p_4\gamma_\nu
 \right].
 \end{align*}
使用 \eqnrefs{eq:gamma-two-sandwich-r10,eq:gamma-three-sandwich-r10} 的 $\gamma$ 矩阵夹心恒等式，
\begin{align*}
 \mathcal I_{st}
 &=-32(p_1\!\cdot p_4)(p_2\!\cdot p_3),\\
 &=-8u^2,
\end{align*}
所以
\begin{equation*}
 2\Re\,\overline{\mathcal M_s\mathcal M_t^*}
 =4g_{\rm em}^4\frac{u^2}{st}.
\end{equation*}
三项相加为
\begin{equation*}
 {
 \overline{|\mathcal M_{\rm Bhabha}|^2}
 =2g_{\rm em}^4\left[
 \frac{u^2+t^2}{s^2}
 +\frac{u^2+s^2}{t^2}
 +\frac{2u^2}{st}
 \right].}
 \end{equation*}
代入质心系的 $t,u$ 后
\begin{equation*}
 {
 \frac{\dd\sigma_{\rm Bhabha}}{\dd\Omega}
 =\frac{\hbar^2\alpha^2}{4s}
 \frac{(3+\cos^2\theta)^2}{(1-\cos\theta)^2}.}
 \end{equation*}
\end{derivation}

\begin{derivation}[从\eqnrefs{eq:2to2-cross-section-master-si-r9}到\eqnrefs{eq:2to2-dsigma-dt-si-r68}]
在质心系写

\begin{equation*}
 p_1^\mu=(E_1/c,\mathbf p_i),
 \qquad
 p_3^\mu=(E_3/c,\mathbf p_f),
\end{equation*}
并令 $\theta$ 为 $\mathbf p_i$ 与 $\mathbf p_f$ 的夹角。由 $t=(p_1-p_3)^2$，
\begin{align*}
 t
 &=p_1^2+p_3^2-2p_1\!\cdot p_3,\\
 &=(m_1^2+m_3^2)c^2
 -\frac{2E_1E_3}{c^2}
 +2|\mathbf p_i||\mathbf p_f|\cos\theta.
 \end{align*}
固定 $s$ 时，$E_1,E_3,|\mathbf p_i|,|\mathbf p_f|$ 都不随 $\theta$ 改变，因此
\begin{equation*}
 \frac{\dd t}{\dd\cos\theta}
 =2|\mathbf p_i||\mathbf p_f|.
\end{equation*}
轴对称两体散射有
\begin{equation*}
 \dd\Omega=2\pi\,\dd(\cos\theta),
 \qquad
 \frac{\dd\Omega}{\dd t}
 =\frac{\pi}{|\mathbf p_i||\mathbf p_f|}.
\end{equation*}
代入二体截面主式，
\begin{align*}
 \frac{\dd\sigma}{\dd t}
 &=\frac{\dd\sigma}{\dd\Omega}
 \frac{\dd\Omega}{\dd t},\\
 &=\frac{\hbar^2}{64\pi^2s}
 \frac{|\mathbf p_f|}{|\mathbf p_i|}
 \overline{|\mathcal M|^2}
 \frac{\pi}{|\mathbf p_i||\mathbf p_f|},\\
 &=\frac{\hbar^2\overline{|\mathcal M|^2}}
 {64\pi s|\mathbf p_i|^2}.
\end{align*}
再由 \eqnrefs{eq:cm-momenta-kallen-si-r9}
\begin{equation*}
 |\mathbf p_i|^2
 =\frac{\Lambda_K(s,m_1^2c^2,m_2^2c^2)}{4s},
\end{equation*}
得到
\begin{equation*}
 {
 \frac{\dd\sigma}{\dd t}
 =\frac{\hbar^2\overline{|\mathcal M|^2}}
 {16\pi\Lambda_K(s,m_1^2c^2,m_2^2c^2)}.}
 \end{equation*}
由于 $\Lambda_K$ 具有动量四次方量纲，右端的量纲是面积除以动量平方，与 $\dd\sigma/\dd t$ 一致。
\end{derivation}

\begin{derivation}[从\eqnrefs{eq:ee-mumu-M}到\eqnrefs{eq:mumu-total-massive-si-r10}]

正文\eqnrefs{eq:ee-mumu-M} 的非极化平方振幅由初、末态 Dirac 张量的缩并给出；二体相空间把它转成角分布，角积分再给总截面。考虑 $s\gg \me^2c^2$ 而允许 $s$ 接近 $4m_\mu^2c^2$，因此忽略初态电子质量而保留末态 $m_\mu$。

末态与初态的 Dirac 张量分别为
\begin{align*}
 L_{\mu}^{\rho\sigma}
 &=\Tr[(\slashed p_3+m_\mu c)\gamma^\rho
 (\slashed p_4-m_\mu c)\gamma^\sigma]\\
 &=4\left[
 p_3^\rho p_4^\sigma+p_3^\sigma p_4^\rho
 -(p_3\!\cdot p_4+m_\mu^2c^2)\eta^{\rho\sigma}
 \right],\\
 L_\ee^{\rho\sigma}
 &=4\left(
 p_2^\rho p_1^\sigma+p_2^\sigma p_1^\rho
 -p_1\!\cdot p_2\,\eta^{\rho\sigma}
 \right).
\end{align*}
因此
\begin{align*}
 \overline{|\mathcal M|^2}
 &=\frac{g_{\rm em}^4}{4s^2}L_{e,\rho\sigma}L_\mu^{\rho\sigma}\\
 &=\frac{8g_{\rm em}^4}{s^2}
 \left[
 (p_2\!\cdot p_3)(p_1\!\cdot p_4)
 +(p_2\!\cdot p_4)(p_1\!\cdot p_3)
 +(p_1\!\cdot p_2)m_\mu^2c^2
 \right].
\end{align*}
在质心系中，正文\eqnrefs{eq:mumu-beta-si-r10} 给出
\begin{align*}
 p_1\!\cdot p_2&=\frac{s}{2},\\
 p_1\!\cdot p_3&=\frac{s}{4}
 (1-\beta_{\rm rel}^{(\mu)}\cos\theta),\\
 p_1\!\cdot p_4&=\frac{s}{4}
 (1+\beta_{\rm rel}^{(\mu)}\cos\theta),
\end{align*}
其余两项由 $1\leftrightarrow2$ 对称性得到。再用
$4m_\mu^2c^2/s=1-(\beta_{\rm rel}^{(\mu)})^2$，整理为正文\eqnrefs{eq:mumu-M2-massive-angle-si-r10}。

二体截面母式给
\begin{equation*}
 \frac{\dd\sigma}{\dd\Omega}
 =\frac{\hbar^2}{64\pi^2s}
 \frac{|\mathbf p_f|}{|\mathbf p_i|}
 \overline{|\mathcal M|^2}.
\end{equation*}
由 $|\mathbf p_f|/|\mathbf p_i|=\beta_{\rm rel}^{(\mu)}$ 与
$g_{\rm em}^2=4\pi\alpha$，得到正文\eqnrefs{eq:mumu-dsigma-massive-si-r10}。

最后令 $x=\cos\theta$，$\dd\Omega=2\pi\dd x$，则
\begin{align*}
 \sigma
 &=\frac{\hbar^2\alpha^2\beta_{\rm rel}^{(\mu)}}{4s}
 2\pi\int_{-1}^{1}\dd x\,
 \left[2-(\beta_{\rm rel}^{(\mu)})^2
 +(\beta_{\rm rel}^{(\mu)})^2x^2\right]\\
 &=\frac{\pi\hbar^2\alpha^2\beta_{\rm rel}^{(\mu)}}{s}
 \left(2-\frac{2}{3}(\beta_{\rm rel}^{(\mu)})^2\right).
\end{align*}
代入
$(\beta_{\rm rel}^{(\mu)})^2=1-4m_\mu^2c^2/s$，得到正文\eqnrefs{eq:mumu-total-massive-si-r10}。
\end{derivation}


